\section{Summary：Lecture 13: Data I 的核心主线}

前面的课程默认“给定数据”，然后讨论如何训练语言模型。Lecture 13 开始反过来问：\textbf{到底训练什么数据？} 课程的核心结论很直接：data does not fall from the sky。训练数据不是“整个互联网”的自然副产品，而是一条从 live service 到 raw data，再到 processed data 的工程、法律和伦理管线。

\begin{importantbox}{本讲总结先行}
\begin{itemize}
    \item Live service \(\rightarrow\) raw data \(\rightarrow\) processed data，中间有 crawling、HTML extraction、filtering、deduplication、PII handling、license checks。
    \item 数据是区分语言模型的重要因素，因为架构和训练流程越来越公开，数据细节却常被保密。
    \item 数据管线涉及 copyright、privacy、robots.txt、ToS、licensing、fair use、shadow libraries 等问题。
    \item 大部分 pipeline 是启发式的，仍有大量改进空间。
\end{itemize}
\end{importantbox}

\begin{knowledgebox}{课堂提示：Data does not fall from the sky}
课程开场先把数据工作从“下载语料”改写成一条生产链：live service 只是不断变化的在线系统，raw data 是抓取或导出的原始对象，processed data 才是经过转换、过滤与去重后可进入训练的材料。老师把这句话放在 Summary 中，是为了提醒读者：数据不会自动出现，真正困难的是持续获取、判断与维护。
\end{knowledgebox}

